\section*{ВВЕДЕНИЕ}
\addcontentsline{toc}{section}{ВВЕДЕНИЕ}

Современные системы беспроводной связи, такие как сотовые сети 4G/5G, Wi-Fi и системы интернета вещей, развиваются в условиях жестких ограничений на выделяемую полосу частот и излучаемую мощность. Для обеспечения высокой скорости передачи данных и надежности соединения необходимо детально понимать процессы, происходящие при распространении радиосигналов в пространственном канале.

Одной из главных проблем беспроводной связи является многолучевое распространение сигналов, вызванное отражениями от зданий, рельефа местности и других препятствий. Это приводит к интерференции, замираниям (федингам) и межсимвольным искажениям. Для борьбы с этими явлениями в приемных устройствах используются методы оценивания канала связи и эквализации (оценивания информационных символов). Качество этих процедур напрямую определяет вероятность битовой ошибки и пропускную способность системы.

Целью данной работы является изучение статистических свойств беспроводных каналов связи и анализ подходов к оцениванию его импульсной характеристики.

В ходе работы решались следующие задачи:
\begin{enumerate}
    \item Рассмотреть базовые модели распространения сигналов (свободное пространство, влияние земной поверхности).
    \item Изучить природу крупномасштабных и мелкомасштабных замираний сигнала.
    \item Проанализировать понятия импульсной характеристики и передаточной функции многолучевого канала, а также временной дисперсии.
    \item Рассмотреть базовый алгоритм оценивания импульсной характеристики канала при известной ее длине.
\end{enumerate}

В качестве основного метода исследования используется аналитический обзор и анализ теоретических основ обработки сигналов в беспроводных системах, представленных в научно-технической литературе.